\chapter{\IfLanguageName{dutch}{Samenvatting}{Abstract}}%
\label{ch:samenvatting}
Via de OBD-II-interface kunnen verschillende gegevens uit een voertuig worden
uitgelezen, zoals motortoerental, snelheid en temperatuurwaarden. Deze gegevens
kunnen gebruikt worden om de toestand van een voertuig realtime op te volgen.
Het uitlezen van de gegevens vormt echter slechts een deel van een
monitoringsysteem. De data moet ook verwerkt, doorgestuurd en op een
overzichtelijke manier aan de gebruiker weergegeven worden.

Deze bachelorproef onderzoekt daarom hoe OBD-II-voertuigdata realtime kan worden
uitgelezen, verwerkt en gevisualiseerd in een webgebaseerde
monitoringapplicatie. Daarbij wordt specifiek onderzocht welke van de
communicatiemethoden REST API en WebSockets het meest geschikt is voor het
doorsturen van voertuigdata tussen de backend en frontend.

Voor het onderzoek werd voertuigdata verzameld van meerdere voertuigen met
behulp van een Vgate VLinker FS OBD-II-adapter en de Pythonbibliotheek
\texttt{python-OBD}. Hierbij werd onderzocht welke voertuigparameters
beschikbaar zijn en welke verschillen tussen voertuigen voorkomen. Vervolgens
werd een monitoringapplicatie ontwikkeld met een FastAPI-backend en een
React-frontend. Daarnaast werd een simulatieomgeving ontwikkeld waarmee de
applicatie ook zonder een fysiek voertuig kan worden ontwikkeld en getest.

REST API en WebSockets werden binnen dezelfde applicatie geïmplementeerd en
onder vergelijkbare omstandigheden getest. Hiervoor werden update-intervallen
van 100, 250 en 1000 milliseconden gebruikt. Iedere configuratie werd gedurende
60 seconden getest en meerdere keren herhaald. De vergelijking richt zich
voornamelijk op de gemiddelde latency, P95-latency en het netwerkverbruik.

% ============================================================================
% TODO NA METINGEN
% Hier de belangrijkste resultaten van de vergelijking invullen.
% ============================================================================

Uit de resultaten bleek dat [REST/WebSockets] binnen de onderzochte
testopstelling gemiddeld de laagste latency behaalde. Ook op het vlak van
P95-latency bleek [RESULTAAT]. Wat het netwerkverbruik betreft, veroorzaakte
[REST/WebSockets] [MINDER/MEER] netwerkverkeer. Een update-interval van
[100/250/1000] milliseconden bleek uiteindelijk de meest geschikte balans te
bieden tussen de actualiteit van de voertuigdata en de benodigde communicatie.

% EINDE TODO

Het onderzoek toont aan dat OBD-II kan worden gebruikt als basis voor een
webgebaseerde realtime monitoringapplicatie. Daarnaast blijkt dat de keuze van
de communicatiemethode en de updatefrequentie invloed heeft op de manier waarop
voertuigdata naar de frontend wordt doorgestuurd. De resultaten moeten worden
geïnterpreteerd binnen de gebruikte testopstelling en het beperkte aantal
onderzochte voertuigen.